\section{Internal slice solver and balanced profiles}
\label{sec:internal-solver-and-profiles-for-these-tilings}

Section~\ref{sec:patch-tiling-under-9} fixed the internal cube slices and
then supplied cleaned first supports for prescribed comparison laws. This section first
constructs the internal odd and even laws for a cluster patch, then chooses
the comparison laws consistently through a fixed point.

A parameter point consists of one input law \(\pi_i\) and one nonnegative price vector \(z_i\) for each cluster patch, with
\[
 M_i\max\pi_i\le e^{10k_iT_i},\qquad
 z_i\in\Delta(Y_i):=\{z\ge0:\ \sum_{y\in Y_i}z(y)=1\}.
\]
This is a compact polytope. Triangulate it finely enough for Lemma~\ref{lem:stable-patch-cleaning}, and so that coordinate price changes within a simplex are \(o(1/M_i)\). At every vertex of this mesh, store the cleaned first laws and the price masks defined below. A fixed parameter point supplies its barycentric probabilities on the vertices of its simplex. Each center and each odd group will draw a vertex independently with these same probabilities; they do not share one vertex draw. The output laws and the eventual fixed point are constructed after the slice experiment.

\begin{proposition}[The internal slice solver]\label{prop:internal-slice-solver}
Fix a cluster patch from Proposition~\ref{prop:patch-extraction-allocation}, cleaned by Lemma~\ref{lem:stable-patch-cleaning} at the corners of a sufficiently fine parameter mesh. For any common choice of barycentric corner probabilities, there is a slice experiment with primitive data \(W\), odd bin probabilities \(q_g^W\), conditional in-bin laws \(U_g^W\), odd marginal laws \(p_g^W\), and even subprobability rows \(\sigma_v\). The odd bin probabilities, individual label masses, and in-bin support sizes obey \eqref{eq:source-13}. The even rows have the bounds
\[
 N\max_x\sigma_v(x)\le 2^{h_i}e^{-500s_i},\qquad
 \mathbb E\sigma_v(x)\le K/M_i,
\]
where the expectation uses unrestricted primitive data and internal reference sampling. Whenever an even row is nonzero, it is a probability law supported in one cleaned corner support and on common neighbors of every internal odd label. The local tests in \eqref{eq:source-15} assert local validity and bound the internal reference probability of \(\sigma_v=0\), conditional on \(W\), by \(\epsilon_i=e^{-.001a k_i h_i}\). They hold simultaneously in a slice except with probability \(\exp(-h_i^{1+\Omega(1)})\). All estimates are uniform in the corner probabilities, with the deterministic local consultation domains specified below.
\end{proposition}

\begin{proof}\leavevmode\par
\subsection{Primitive data and selection geometry}
\label{sec:internal-primitive-data}

Fix an \(h_i\)-slice and abbreviate \(h=h_i,k=k_i,T=T_i\). Conditional on the common parameter point, different slices have independent copies of the following experiment. At every potential center \(c\), indexed by an internal location and a level, draw a mesh vertex, use its cleaned law \(\mu_c\), and draw \(W_c\sim\mu_c^{\otimes k}\). These variables are generated even when the center is absent. Each odd group independently draws its mask vertex. Prospective positions, the orders used in searches, activations, and site-level choice ties are further independent primitive variables. Eligibility is computed before using activations or choice ties. We write \(W\) for this whole pre-bin history, including the selections derived from the primitives. Its derived fields need not be independent. Before success conditioning, the records at distinct centers and groups and the remaining randomizers are independent; each tuple is sampled conditional on its own corner.

Here is the projection defining groups and sites. Index the \(h\) internal coordinates by the binary group of order \(h\), and for a word \(z\) set
\[
 s(z)=\sum_l lz_l,\qquad \psi(z)=z\oplus e_{s(z)}.
\]
The map has syndrome-zero image and fibers of size \(h\), and always flips one bit. An odd group is a fiber on the odd roles of the slice; a selection site is \(\psi(v)\) for an even role. The parity convention includes the shift from the slice's fixed complementary word. A group's site neighborhood consists of the projected even roles adjacent to one of its odd roles. As in Section~\ref{sec:projections-and-lists}, these sites have mutual distances at most 6, and each site lies in \(O(h^3)\) such neighborhoods.

For a fixed even role,
\[
 \psi(v\oplus e_l)=v\oplus e_l\oplus e_{s(v)+l}.
\]
If \(s(v)\ne0\), the flips \(l\) and \(l+s(v)\) form pairs with the same image; if \(s(v)=0\), all \(h\) flips have the same image. Thus each group queried by this star has multiplicity \(j\ge2\), and the multiplicities sum to \(h\). This is the incidence property used in the posterior calculation.

For each group and mesh vertex, form a mask from that vertex's price vector on \(Y_i\). Take the uniform prior on bins with uniform in-bin laws; retain bins having at least half their labels of price \(\le10/M_i\), then restrict to those cheap labels in each bin, normalizing both stages. At most a tenth of labels overall exceed the cut, so most bins have enough cheap mass; aggregate is bounded by 4 times uniform on \(Y_i\). A physical bin is a member of the fixed partition of \(Y_i\). Write \(D\) for its masked uniform probability law when it occurs inside a formula such as \(D(F)\), and \(r_g(D)\) for the masked prior probability of choosing that bin at group \(g\). Thus a bin identity determines a group-dependent masked law. All these prior lookups and draws up to the actual assignment of clusters are local. In particular different centers can have different cleaned first laws, all still giving codegrees \(\ge1/4+a\).

We give details of the template, including its gates used later. Use the geometric height device on \(Q_h\) with step bound 6, radius \(\lfloor\rho h\rfloor\), \(\lambda=h^{10}\), activation and crowd exponents \(b_0=\omega/8,\ b=\omega/2\), height-proof exponents \(a_g=\omega/12,\ \zeta_g=1-\theta_g=\omega/30,\ 0<\sigma_g\ll\omega\) (subscripts refer here to geometric parameters only). Positions and activations are independent of tuples, masks and tie randomizers. The consultation distance for paths is \(o(h)\) in this geometric estimate, in particular tiny compared to \(\rho h\) by the large threshold. Throughout, asymptotic slack in powers of \(h\) itself need only hold for \(q\ge Q_0\), choosing that fixed threshold sufficiently large.

\subsection{List tests and odd bin laws}
\label{sec:internal-list-tests}

Fix corners, the group's mask, prospective positions, and a hypothetical set \(\mathcal D\) of 1 through \(T\) distinct candidate IDs in its neighborhood. Let \(F\) be the set hitting all entries of their tuples and let \(F_{-c}\) omit tuple \(c\). Set \(A(J)=\mathbb E_{D\sim r_g}D(J)^2\). The list must satisfy the absolute and deletion bounds from \eqref{eq:source-4}, with every ID treated as an own-slice ID:
\[
 A(F)\ge e^{-2k|\mathcal D|},\qquad
 \frac{A(F)}{A(F_{-c})}\ge e^{(-\log4+.4a)k}
 \quad(c\in\mathcal D).
\]

To bound failure, expose tuple coordinates successively; for each deletion test choose an order in which the indicated tuple is last. If the entries already exposed impose the hit set \(J\), the next ratio is
\[
 \frac{A(J\cap N_G(x))}{A(J)}
   =\mathbb E_{D\text{ tilted by }D(J)^2}d_G(x;D|_J)^2.
\]
The next \(x\) has its center's cleaned first law. All-pairs codegrees give conditional mean ratio at least \(1/4+a\), even though different centers can use different corner laws. While preceding ratios are at least \(.24\), the second aggregate law has width at most
\(\log(N/M_i)+O(1)+2kT<n^{\eta_0}/2\).
Jensen and \eqref{eq:source-2} then bound the probability of a ratio below \(1/4-O(n^{-\eta_0})\) by \(e^{-\Omega(n^{\eta_0})}\).

As in the proof of \eqref{eq:source-4}, clip log ratios below at \(\log(.24)\). Their conditional means are at least \(-\log4+.7a\). The deletion inequality asks the last \(k\) increments to sum to at least \((-\log4+.4a)k\), leaving a deviation margin \(.3ak\). Bounded-increment concentration gives failure \(e^{-\Omega(a^2k)}\). For the absolute inequality, sum all \(k|\mathcal D|\) increments and use its still larger margin above \(-2k|\mathcal D|\). Stop on an earlier invalid ratio, continue the clipped process with good values, and separately add the negligible probability of an unclipped catastrophe. The finite collection of orders gives the stated failure bound for one list.

In each group neighborhood search failed lists using present IDs in range of incident even sites, at any levels, and mark a maximal disjoint family's union forbidden for choices by those incident sites. On polynomial position bounds there are at most \(\exp(O(T\log h))\) lists per search and disjoint tests in one search are independent conditional on positions, masks and corner realizations. Here disjoint lists have disjoint ID sets. A union bound for \(h\) disjoint failures, using \(T\log h=o(a^2k)\), shows that each search has fewer than \(h\) of them except with probability \(\exp(-h^{1+\Omega(1)})\) per slice. Such a search marks fewer than \(hT\) IDs. Since a site participates in \(O(h^3)\) searches, it loses at most \(O(h^4T)=o(\lambda)\) choices. Binomial position concentration puts the original counts close to \(\lambda\); hence the remaining eligible count is at least \(\lambda/2\), while the original count is at most \(2\lambda\), with the same slice failure bound. Now apply the activation/path mechanism with these eligibilities. Choose one ID per site at the obtained height by a uniform active eligible tie (that same choice used everywhere it is queried). The largest scale proof gives heights and non-bad-layer choices everywhere there except with probability \(\exp(-h^{1+\Omega(1)})\) on size success. Realized lists at each group then have size \(\le T\) by the crowd bound at the (at most two) levels used, and pass \eqref{eq:source-4}: a failed realized list uses no marked ID and would enlarge the maximal disjoint family.

At each group on validity, given \(W\), sample independently a masked \(D\) tilted by \(D(F)^2\), conditional further on
\[
 D(F)\ge e^{-1.5k|\mathcal D|},\qquad
 D(F)/D(F_{-c})\ge e^{(-\log2+.08a)k}\quad\text{for all }c .
\]
The extra restrictions retain probability at least \(1/2\) by \eqref{eq:source-4}. Thus the first stage selects a physical bin by a restricted squared-mass tilt of \(r_g\). Conditional on that bin, the internal reference law draws all queried role labels independently from its uniform law \(D|_F\). Write \(q_g^W(D)\) for the induced bin probabilities at group \(g\), \(U_g^W(D)\) for the resulting individual conditional uniform-subset law given bin, \(p_g^W=\sum_D q_g^W(D) U_g^W(D)\). On needed fallbacks just draw from the group's masked prior and laws. The label-marginal cap uses cancellation of one hit-mass factor:
\[
 p_g^W(y)\le\frac{2}{A(F)}\sum_D r_g(D)D(F)D(y)\mathbf1_F(y)
       \le\frac{8e^{2kT}}{M_i}.
\]
The masked bin prior is at most \(2d_i/M_i\), so its squared tilt and extra restrictions give the first cap below. The absolute mass restriction and the masked support size give the third. These bounds also hold for the prescribed fallbacks:
\begin{equation}
 q_g^W(D)\le 4e^{2kT} d_i/M_i,\quad p_g^W(y)\le 8e^{2kT}/M_i,\quad
 |\operatorname{supp} U_g^W(D)|\ge \tfrac12 d_i e^{-1.5kT}
 \label{eq:source-13}
\end{equation}
for used bins. In this notation \(g\) as subscript is a group, not bias budget.

\paragraph{Local validity and fallback rules.}
The local rules are defined even when some tests fail. At a site, use the height rule restricted to its long consultation domain. Choose no ID if that height reaches the horizon or its site-level is bad; otherwise choose a uniform active eligible ID, with the pre-randomized tie for that level.

A group is valid when all its incident sites make choices, every candidate ball at those sites and levels has at most \(2\lambda\) positions, and the resulting list has size between 1 and \(T\) and passes \eqref{eq:source-4}. An invalid group uses the masked-prior fallback above.

For an even star at \(v\), let \(\mathcal V_v\) require validity of its neighboring groups and the position upper bound \(2\lambda\) and eligibility lower bound \(\lambda/2\) at every site-level in its projected site's long consultation domain. This event depends on \(W\), before bins or labels are drawn. It ensures a genuine selection at the projected site. A positive selected level then requires a path reaching positive height, and selection uses the actual independent uniform active tie. All these events hold simultaneously on the slice success already proved.

Fallback/mark/choice tie rules can use separate independent order permutations per queried group or site-level, as needed (on ID lists or on IDs, respectively) to be translation-invariant in law under syndrome-zero shifts, including shifts interchanging the two local parity conventions. Choice ties themselves are independent of eligibility and activation data. All validity and comparison rules may and will use this symmetry. The height lemma with ambient parameter \(h\) gives \(2DH+O(\rho h+1)<10\rho h\) after increasing the threshold. Thus the rules consult only this slice within \(10\rho h\) of the inner queried state, including in the counterfactual calculations below (eligibility calculations themselves just use incident group tests in range, not their selections).

\subsection{Even posteriors and predictive tests}
\label{sec:internal-even-posteriors}

Fix a present candidate ID \(c\) for \(v\)'s projected site. Hold all other primitive data fixed, including the corner selecting its prior \(\mu_c\), and vary \(w=W_c\). For ordered internal neighbor labels \(\boldsymbol y\), define
\[
 L_w(\boldsymbol y)=
 \mathbf1_{\{c\text{ selected},\,\mathcal V_v\}}(W_{-c},w)\,
 \Pr_{\rm int}\{\boldsymbol Y=\boldsymbol y\mid W_{-c},w\}.
\]
The bin at each group is integrated out in this probability. Selections, validity, and all bin laws are recomputed at \(w\).

We construct a reference independent of \(w\). For each queried group \(g\), let \(\mathcal L_g\) be the family of all lists of at most \(T\) present IDs in its candidate ranges that contain \(c\). This family uses only the frozen positions. For \(\mathcal D\in\mathcal L_g\), delete \(c\), tilt the masked bin prior by \(D(F_{-c})^2\), and draw its \(j_g\) queried labels independently from \(D|_{F_{-c}}\). Denote this probability by \(Q_{g,\mathcal D}\), using a fixed symmetric fallback if its normalization is zero. Average uniformly over the family, with a fallback also for an empty family, and set
\[
 \mathcal Q=\bigotimes_g
  \left(\frac1{|\mathcal L_g|}\sum_{\mathcal D\in\mathcal L_g}Q_{g,\mathcal D}\right).
\]
This notation uses the fallback convention when the family is empty. In particular the actual list may depend on \(w\), but the reference mixture does not.

For one fixed list, the likelihood ratio before integrating its bin is at most
\[
 2\frac{A(F_{-c})}{A(F)}
       \left(\frac{D(F)}{D(F_{-c})}\right)^{2-j_g}
 \le2e^{(\log2-.08a)kj_g}.
\]
The factor 2 pays for the extra bin restrictions; the other factors are the squared-tilt normalizer and the \(j_g\) conditional label densities. The inequality uses \(j_g\ge2\) and the tested ratio bounds. Passing to the fixed mixture costs \(|\mathcal L_g|\), with logarithm \(O(T\log h)=o(ak)\) on the position gate. Since \(\sum_g j_g=h\) and there are no singleton incidences,
\[
 L_w\le e^{(\log2-.06a)kh}\mathcal Q.
\]
If a required position bound fails, the sublikelihood is zero.

With \(c\) actually selected on \(\mathcal V_v\), require
\[
 Z(\boldsymbol y)=\int L_w(\boldsymbol y)\,d\mu_c^{\otimes k}(w)>0,
 \qquad Z(\boldsymbol y)\ge e^{-.01akh}\mathcal Q(\boldsymbol y).
\]
The posterior \(P_{c,\boldsymbol y}(dw)=L_w(\boldsymbol y)d\mu_c^{\otimes k}(w)/Z(\boldsymbol y)\) has density at most \(e^{(\log2-.05a)kh}\) relative to its prior. Since \(\max\mu_c\le2/M_i\) and the allocation gives \(\log(2N/M_i)\le .01ah\), its joint atom cap relative to \(N^{-k}\) is \(e^{(\log2-.04a)kh}\).

Let \(\bar P\) be its average coordinate marginal, and let
\(B=\{x:N\bar P(x)>2^he^{-.02ah}\}\). Then \(|B|/N\le2^{-h}e^{.02ah}\). For any specified \(r=\lceil(1-.01a)k\rceil\) coordinates, the joint cap bounds the probability that all lie in \(B\) by
\[
 e^{(\log2-.04a)kh}(|B|/N)^r=e^{-\Omega(akh)}.
\]
The union over at most \(2^k\) coordinate subsets is negligible. Hence the expected fraction of coordinates outside \(B\), namely \(\bar P(B^c)\), is at least \(c a\) for a fixed \(c>0\). Define \(\sigma_v\) by restricting \(\bar P\) to \(B^c\) and normalizing; this costs \(O(1/a)\), which the fixed threshold absorbs in the cap below.

On failures of the gates let \(\sigma_v=0\); otherwise it is supported in a single cleaned corner law's support and on common neighbors of \textbf{all} internal labels. We have
\begin{equation}
 N\max_x\sigma_v(x)\le 2^h e^{-500s_i},\qquad
 \mathbb E\sigma_v(x)\le K/M_i
 \label{eq:source-14}
\end{equation}
where expectation is for unrestricted \(W\) and internal reference sampling and \(K\) a fixed constant. To prove the mean bound, first fix the data other than \(W_c\). If \(\bar P_{c,\boldsymbol y}\) is the posterior coordinate marginal, cancellation of its denominator gives
\[
 \sum_{\boldsymbol y} Z(\boldsymbol y)\bar P_{c,\boldsymbol y}(x)
 =\int \frac1k\sum_{r\le k}\mathbf1_{\{w_r=x\}}
       \sum_{\boldsymbol y}L_w(\boldsymbol y)\,d\mu_c^{\otimes k}(w).
\]
The inner sum is the probability of selecting \(c\) and passing \(\mathcal V_v\) in this candidate experiment. Thus, after the \(O(1/a)\) restriction cost, the contribution is bounded by the actual tuple-coordinate incidence with that selection gate.

An entry equals \(x\) with prior probability at most \(2/M_i\); conditioning on this equality leaves positions and activations with their original laws. Given presence of \(c\), the gated selection probability at level zero is \(O(1/\lambda)\) by the uniform active tie bound. At higher levels the position-averaged, eligibility-uniform positive-height bound is \(e^{-h^{\Omega(1)}}\), with the one specified center forced present. These are precisely the bounds of Lemma~\ref{lem:height-selection}; they do not condition on global success. The sum of candidate presence probabilities is \(\lambda\) per level. Summing over the polynomial number of levels proves \(\mathbb E\sigma_v(x)\le K/M_i\).

Let \(E_v^{\rm in}\) denote failure (\(\sigma_v=0\)). For each candidate, Lemma~\ref{lem:gated-posterior} bounds failure of the displayed predictive inequality, on its selection and validity gate, by \(e^{-.01akh}\). The position-count gate permits only polynomially many candidates. Consequently the nonnegative function of primitive data
\[
 B_v(W):=\mathbf1_{\mathcal V_v}\Pr_{\rm int}(E_v^{\rm in}\mid W)
\]
has raw mean at most \(h^{O(1)}e^{-.01akh}\). For \(\epsilon_i=e^{-.001a k_i h_i}\),
\begin{equation}
 \mathcal H_v:\quad \mathcal V_v,\qquad
     \Pr_{\rm int}(E_v^{\rm in}\mid W)\le\epsilon_i
 \label{eq:source-15}
\end{equation}
holds simultaneously per slice except with probability \(\exp(-h_i^{1+\Omega(1)})\): apply Markov to \(B_v\) at threshold \(\epsilon_i\), union over the \(2^{h-1}\) even roles, and add the separate slice-wide validity failure. The exponent \(akh\), with \(k=h^{3\omega}+O(1)\), pays for this union. The individual tests \eqref{eq:source-15} are local as well. All these estimates are uniform in the corner probabilities.
\end{proof}

\subsection{The simultaneous profile choice}
\label{sec:internal-profile-choice}

Use the parameter polytope and mesh defined at the start of this section. For each cluster patch, define its unrestricted mean odd law by
\[
 \pi_i^0=\mathbb E_W p_g^W.
\]
In high mode set \(\pi_i^{\rm out}=\pi_i^0\). In low mode first condition the primitive slice history \(W\) on \eqref{eq:source-15} at every even role. For such a history, let \(\mathcal B_g(W)\) be the positive bins satisfying
\[
 \Pr_{\rm int}(E_v^{\rm in}\mid W,D_g=D)\le\sqrt{\epsilon_i}
 \quad\text{for every internal star incident to }g.
\]
There are at most \(h^2\) incident stars. Markov and \eqref{eq:source-15} give \(q_g^W(\mathcal B_g(W))\ge1-h^2\sqrt{\epsilon_i}\), so the pretrimmed bin law
\[
 q_g^{\mathrm{in},W}(D)=
 \frac{q_g^W(D)\mathbf1_{\mathcal B_g(W)}(D)}
      {q_g^W(\mathcal B_g(W))}
\]
is well defined. Leave the in-bin laws \(U_g^W(D)\) unchanged, and define the low-mode output \(\pi_i^{\rm out}\) as the mean of \(\sum_Dq_g^{\mathrm{in},W}(D)U_g^W(D)\) under the conditioned slice law. The tests here evaluate the original raw failure rule; no posterior is recomputed with a trimmed prior. Syndrome-zero translation symmetry makes these means independent of the group \(g\), including across the two parity conventions.

\begin{proposition}[Balanced patch profiles]\label{prop:balanced-patch-profiles}
In the capped product domain of cluster-patch input laws and price simplices, there is a parameter point \((\pi,z)\) such that \(\pi_i=\pi_i^{\rm out}\) and \(z_i\) maximizes \(z_i'\cdot\pi_i\) over \(\Delta(Y_i)\), for every patch \(i\). At this point,
\begin{equation}
 \max\pi_i\le 11/M_i . \label{eq:source-16}
\end{equation}
Every cleaned corner support used with positive probability satisfies Lemma~\ref{lem:stable-patch-cleaning} at the realized comparison laws \((\pi_i)_i\). Their union in a patch, called its envelope, inherits the degree and row-tail bounds; clique exclusion is asserted separately on each corner support. In low cluster mode the unrestricted mean satisfies
\[
 \|\pi_i^0-\pi_i\|_{\rm TV}
 \le \exp(-h_i^{1+\Omega(1)})+O(h_i^2\sqrt{\epsilon_i}).
\]
\end{proposition}

\begin{proof}
The output means are continuous in the parameters. Barycentric probabilities are continuous and piecewise affine; varying them changes the weights of the finite primitive experiments but leaves every corner lookup, recomputed validity gate and posterior rule for a fixed realization unchanged. In particular, each posterior retains its realized corner law as prior. The slice-conditioning denominators are bounded below, as are the pretrim normalizers just estimated. Equation~\eqref{eq:source-13} puts the outputs in the capped input domain: the bounded conditioning and trim factors fit the gap from \(2k_iT_i\) to \(10k_iT_i\).

Apply Kakutani to the response
\[
 (\pi,z)\longmapsto
 \left(\pi^{\rm out}(\pi,z),\,
       \prod_i\arg\max_{z_i'\in\Delta(Y_i)}z_i'\cdot\pi_i\right).
\]
This response is nonempty, convex and has closed graph, so it has a fixed point. There \(z_i\cdot\pi_i=\max_y\pi_i(y)\). Every label in the output support is cheap at some mesh vertex used by the point, whose price differs coordinatewise from \(z_i\) by \(o(1/M_i)\). Hence \(z_i(y)\le11/M_i\) on that support. Integrating against \(\pi_i=\pi_i^{\rm out}\) proves \eqref{eq:source-16}.

The mesh was chosen for the stability conclusion of Lemma~\ref{lem:stable-patch-cleaning}. Each used corner support therefore has the cleaned properties at the fixed-point laws. Degree and row-tail bounds hold label by label and pass to the envelope; the no-clique condition remains a property of each corner support separately.

In low mode, conditioning on all slice gates changes the raw history law in total variation by at most their failure probability, \(\exp(-h_i^{1+\Omega(1)})\). For each conditioned history, the bin pretrim changes the label marginal by at most \(h_i^2\sqrt{\epsilon_i}\). Averaging and the triangle inequality give the stated comparison with \(\pi_i^0\). These bounds are uniform above the fixed threshold \(Q_0\); they do not require a low-mode scale to tend to infinity.
\end{proof}

\begin{lemma}[Finite low-mode data]\label{lem:finite-low-mode-data}
At the fixed point, the low-mode primitive slice data can be represented with log cardinality at most \(n^{1+o(1)}\).
\end{lemma}

\begin{proof}
Encode \(W\) by its finite categorical outcomes, not by continuous random seeds. A parameter point uses at most dimension plus one active vertices, hence \(O(N)\) possibilities for each corner draw. In low mode \(h\le(\log n)^{.1}\), so the number \(h^{O(1)}2^h\) of position, group, and level records is \(n^{o(1)}\). Each tuple and vertex record costs at most \(O(k\log N)=n^{1+o(1)}\) in log cardinality, using \(N\le n2^n\).

The order universes for IDs and lists have sizes
\(U=\exp(O(T(h+\log h)))=n^{o(1)}\). An order contributes at most \(\log(U!)\le U\log U=n^{o(1)}\). Multiplying by the number of records, and including the binary presence and activation variables, gives total log cardinality \(n^{1+o(1)}\).
\end{proof}

In direct modes none of the internal mechanism is needed: \(\sigma_v\) always denotes the uniform cleaned first law in the patch and \(\pi_i\) the fixed uniform second law.

The calibrated near-product bound~\eqref{eq:source-17}, from Lemma~\ref{lem:near-product-injection}, will be used for assignments within bins.
